\chapter{KNOWLEDGE AND SKILLS ACQUIRED}

\section{Technical Knowledge}

\subsection{Technical Topic 1}
I learned how an offline-first browser extension can be connected to a cloud
backend without breaking the user flow. The project showed how a local cache,
authentication, REST APIs, and cloud persistence can work together.

\subsection{Technical Topic 2}
I also learned the tradeoffs of serverless architecture. AWS Lambda, API
Gateway, DynamoDB, and S3 reduced operational overhead, but they required
careful attention to permissions, deployment configuration, and request flow.

\subsection{Security, Quality, or Reliability}
The internship reinforced the importance of least-privilege IAM design, CORS
configuration, budget guardrails, log retention, and testable deployment steps.
Reliability also mattered in the offline sync flow, where local cards had to be
preserved until the cloud upload succeeded.

\section{Technical Skills}

\begin{itemize}
  \item Building and debugging a Chrome extension with Manifest V3.
  \item Designing and testing a Node.js REST API backed by DynamoDB.
  \item Deploying and verifying AWS infrastructure with SAM and the AWS CLI.
\end{itemize}

\section{Soft Skills}

\begin{itemize}
  \item Communicating progress and blockers more clearly.
  \item Planning work in weekly milestones and documenting evidence.
  \item Learning unfamiliar tools independently before asking for help.
\end{itemize}

\section{Connection to University Knowledge}

The internship connected directly to software engineering, web development,
databases, operating systems concepts, cloud computing, and security. It also
reinforced the value of testing, version control, architecture planning, and
technical documentation learned at university.